\chapter{Control Flow}
\label{ch:control}

\standfirst{There are no control structures. There are messages that happen to
be sent to booleans, blocks and collections, and the effect is
indistinguishable---except that you can add your own.}

\section{Conditionals}

\ct{ifTrue:} is a message to a Boolean. \ct{True} implements it by evaluating
its argument; \ct{False} implements it by answering \ct{nil}. That is the
entire mechanism.

\begin{st}
x > 0
    ifTrue: ['positive']
    ifFalse: ['not positive']

x isNil ifTrue: [^self error: 'no x'].

result := x > 0 ifTrue: ['yes'].     "nil when x <= 0"
\end{st}

The four forms are \ct{ifTrue:}, \ct{ifFalse:}, \ct{ifTrue:ifFalse:} and
\ct{ifFalse:ifTrue:}. All of them answer the value of whichever block ran, and
\ct{nil} when none did.

\subsection{And, or, and short-circuiting}

\begin{st}
a and: [b]           "short-circuit: b is only evaluated if a is true"
a or: [b]            "short-circuit"
a & b                "both always evaluated"
a | b                "both always evaluated"
a not
\end{st}

The block brackets in \ct{and:} and \ct{or:} are what make them
short-circuiting: the argument is a block precisely so it need not be run. Get
into the habit of writing them, because \ct{a and: b} is legal, evaluates
\ct{b} unconditionally, and will eventually bite you.

\section{Loops}

\subsection{While}

The receiver is a block too, because it has to be re-evaluated each time
round:

\begin{st}
| n |
n := 0.
[n < 5] whileTrue: [n := n + 1].
n                                    "5"
\end{st}

\ct{whileFalse:} is the mirror image. \ct{whileTrue} and \ct{whileFalse}
without arguments loop on the receiver alone. \ct{repeat} loops forever, and
you leave it with a non-local return.

\subsection{Counted}

\begin{st}
1 to: 5 do: [:i | ... ]              "i = 1, 2, 3, 4, 5"
1 to: 10 by: 3 do: [:i | ... ]       "1, 4, 7, 10"
10 to: 1 by: -1 do: [:i | ... ]      "counts down"
5 timesRepeat: [ ... ]               "when you do not need the index"
\end{st}

\subsection{Over a collection}

Almost always what you actually want:

\begin{st}
aCollection do: [:each | ... ]
aCollection do: [:each | ... ] separatedBy: [ ... ]
aCollection reverseDo: [:each | ... ]
aCollection doWithIndex: [:each :i | ... ]
aDictionary keysAndValuesDo: [:k :v | ... ]
\end{st}

Chapter~\ref{ch:collections} has the full set, including the ones that answer
something useful---\ct{collect:}, \ct{select:}, \ct{detect:},
\ct{inject:into:}---which you should reach for before \ct{do:} with an
accumulator.

\section{Why this works, and what it costs}

Sending a message to evaluate a block for every iteration of a loop sounds
expensive. It is not, because the compiler recognises a fixed set of these
patterns and compiles them to jumps, creating no block and sending no message:

\begin{center}
\small
\begin{tabular}{@{}l@{}}
\toprule
\ctp{ifTrue:} \ctp{ifFalse:} \ctp{ifTrue:ifFalse:} \ctp{ifFalse:ifTrue:}\\
\ctp{and:} \ctp{or:}\\
\ctp{whileTrue:} \ctp{whileFalse:} \ctp{whileTrue} \ctp{whileFalse}\\
\ctp{to:do:} \ctp{to:by:do:}\\
\bottomrule
\end{tabular}
\end{center}

Everything else---\ct{do:}, \ct{timesRepeat:}, \ct{detect:}, all of
it---is a genuine message send with a genuine block. That is usually fine, and
where it is not, the fix is a better algorithm rather than a special form.

\begin{stnote}
Inlining is why the receiver of an \ct{ifTrue:} must really be a Boolean. The
jump bytecode tests for \ct{true} and \ct{false} and complains if it finds
anything else, so a nil receiver reports \ct{must be a boolean} rather than
\ct|doesNotUnderstand: #ifTrue:|. When you see that message, the usual cause
is an uninitialised instance variable.
\end{stnote}

\section{Early exit}

There is no \ct{break} and no \ct{continue}. To leave a loop early, return
from the method with \ct|^|:

\begin{st}
indexOf: anObject
    1 to: self size do: [:i |
        (self at: i) = anObject ifTrue: [^i]].
    ^0
\end{st}

To skip an iteration, put the rest of the body inside an \ct{ifTrue:}. If that
makes the block ugly, that is the language suggesting the body wants to be its
own method.

When you need to leave a loop without leaving the method, wrap the loop in a
block that takes an exit and send it \ct{valueWithExit}. Sending \ct{value:}
to the exit leaves the block at once, and \ct{valueWithExit} answers what the
exit was given---or the block's own value when the exit is never taken:

\begin{st}
[:exit |
    #(1 2 3) do: [:each |
        each > 2 ifTrue: [exit value: each]].
    #allGood] valueWithExit                "3"
[:exit |
    #(1 2) do: [:each |
        each > 2 ifTrue: [exit value: each]].
    #allGood] valueWithExit                "#allGood"
\end{st}

\ldots or, more simply, use \ct{detect:ifNone:}, which is that pattern with a
name.

\section{Case statements}

There are none, and the omission is deliberate. The three replacements, in
order of preference:

\begin{enumerate}
\item \textbf{Polymorphism.} Give each case its own class and its own method.
\item \textbf{A Dictionary of blocks.} When the cases are data rather than
  behaviour:
\begin{st}
| handlers |
handlers := Dictionary new.
handlers at: #add put: [:a :b | a + b].
(handlers at: op) value: x value: y
\end{st}
\item \textbf{A chain of \ct{ifTrue:}s.} Honest, readable, and fine for three
  cases.
\end{enumerate}

\section{Writing your own}

Because control flow is messages and blocks, adding to it needs no compiler
change. Here is a complete new control structure:

\begin{st}
Number >> upToBy: stop by: step do: aBlock
    | i |
    i := self.
    [i <= stop] whileTrue: [
        aBlock value: i.
        i := i + step]
\end{st}

And here is one you might actually want---running a block with a guaranteed
teardown, on a receiver that supplies it:

\begin{st}
Object >> withLock: aBlock
    self lock acquire.
    ^aBlock ensure: [self lock release]
\end{st}

That is the whole of what a language designer would have had to add as
syntax. This is the reason Smalltalk's grammar has not needed to grow.
